%% kapa: exam=1 round=19 year=2008 period=1 form=A subject=appraisal_law subject_ko=부동산관계법규 seq=39 number=79
%% source: Q-Net 감정평가사 기출문제 (cst00309) artlSeq=1001962 file=extracted/1008/1001962/2008년 제19회 감정평가사 1차시험 기출문제/1교시.hwp
%% conversion: kordoc 4.24.51 -> tools/gen1.py (auto); review=proofread-sonnet; flags=none
\begin{question}{79}
부동산등기법상 등기의 종류에 관한 설명으로 옳지 않은 것은?
\begin{choices}
  \item 사실의 등기란 부동산의 상황에 관한 등기로서 표제부의 표시란에 하는 등기이다.
  \item 예비등기란 등기원인의 무효 또는 취소로 인한 등기의 말소 또는 회복의 소가 제기된 경우에 제3자를 위해 수소법원의 촉탁으로 행하는 등기를 말한다.
  \item 변경등기란 등기와 실체관계가 후발적으로 일치하지 않는 것을 시정하기 위해 행하는 등기이다.
  \item 경정등기란 등기절차에 착오 또는 유루가 있어서 등기와 실체관계의 원시적 불일치가 생기는 경우에 이를 시정하기 위해 하는 등기이다.
  \item 기입등기란 새로운 등기원인에 의하여 새로운 사항을 등기부에 처음 기입하는 등기를 말한다.
\end{choices}
\end{question}
\answer{2}
